\exercisesection{Exercises for § 1}

\begin{enumerate}
\def\labelenumi{\arabic{enumi})}
\tightlist
\item
  \begin{enumerate}
  \def\labelenumii{\alph{enumii})}
  \tightlist
  \item
    Let (W, S) be a Coxeter system and let \(s_{1}, \ldots, s_{r}\) be elements of S. Put \(w = s_{1} \ldots s_{r}\) . Show that if \(l_{\mathbb{S}}(w) < r\) , there exist two integers p and q with \(1 \leqslant p < q \leqslant r\) such that \(w = s_{1} \ldots s_{p-1}s_{p+1} \ldots s_{q-1}s_{q+1} \ldots s_{r}\) . Show that there is a strictly increasing sequence of integers \(j(1), \ldots, j(k)\) between 1 and r such that \((s_{j(1)}, \ldots, s_{j(k)})\) is a reduced decomposition of w.
  \end{enumerate}
\end{enumerate}

\begin{enumerate}
\def\labelenumi{\alph{enumi})}
\setcounter{enumi}{1}
\tightlist
\item
  Let \((\mathbf{W},\mathbf{S})\) be a Coxeter system and \(X,Y,Z\) three subsets of \(S\) . Show that
\end{enumerate}

\[
W _ {X} \cap (W _ {Y}. W _ {Z}) = (W _ {X} \cap W _ {Y}). (W _ {X} \cap W _ {Z})
\]

(show that every element \(w \in W_{Y}.W_{Z}\) has a reduced decomposition

\[
(s _ {1}, \ldots , s _ {h}, t _ {1}, \ldots , t _ {k})
\]

such that \(s_i \in Y\) and \(t_j \in Z\) and use Cor. 1 of Prop. 7 of no. 8).

Show that

\[
W _ {X}. \left(W _ {Y} \cap W _ {Z}\right) = \left(W _ {X}. W _ {Y}\right) \cap \left(W _ {X}. W _ {Z}\right).
\]

\begin{enumerate}
\def\labelenumi{\arabic{enumi})}
\setcounter{enumi}{1}
\item
  Let \((W,S)\) be a Coxeter system and X a subset of S. Show that \(W_{X}\) is normal in W if and only if every element of X commutes with every element of S-X.
\item
  Let \((\mathbf{W},\mathbf{S})\) be a Coxeter system and \(\mathbf{X},\mathbf{Y}\) two subsets of \(\mathbf{S}\) . Let \(a\in W\) . Show that there exists a unique element \(w\in W_X a W_Y\) of minimal length and that every element
\end{enumerate}

\[
w ^ {\prime} \in W _ {X} a W _ {Y}
\]

can be written in the form \(w' = xwy\) , with \(x \in W_X, y \in W_Y\) and \(l(w') = l(x) + l(w) + l(y)\) (take an element of minimal length in \(W_X aW_Y\) and use Exerc. 1). An element \(w \in W\) is said to be \((X,Y)\) -reduced if it is the element of minimal length in the double coset \(W_X aW_Y\) .

Show that if \(w\) is \((X, \emptyset)\) -reduced then \(l(xw) = l(x) + l(w)\) for all \(x \in W_X\) , and that every element of \(W\) can be written uniquely in the form \(xw\) where \(x \in W_{X}\) and w is \((X, \varnothing)\) -reduced. Show that an element \(w \in W\) is \((X, \varnothing)\) -reduced if and only if \(l(xw) > l(w)\) for all \(x \in X\) (write w in the form \(yw'\) , where \(y \in W_{X}\) and \(w'\) is \((X, \varnothing)\) -reduced).

Show that \(w \in W\) is \((X, Y)\) -reduced if and only if w is both \((X, \varnothing)\) -reduced and \((\varnothing, Y)\) -reduced.

\begin{enumerate}
\def\labelenumi{\arabic{enumi})}
\setcounter{enumi}{3}
\item
  Let \(n\) be an integer \(\geqslant 2\) . For any integer \(i\) with \(1 \leqslant i \leqslant n - 1\) , denote by \(s_i\) the transposition of \(i\) and \(i + 1\) in the set \(\{1, 2, \ldots, n\}\) , and by \(H_i\) the set of \(w \in \mathfrak{S}_n\) such that \(w^{-1}(i) < w^{-1}(i + 1)\) ; put \(S = \{s_1, \ldots, s_{n-1}\}\) . Show that \((\mathfrak{S}_n, S)\) is a Coxeter system and that \(H_i\) is the set of \(w \in \mathfrak{S}_n\) such that \(l(w) < l(s_i w)\) (use Prop. 6 of no. 7).
\item
  Let X be a non-empty set and W a set of permutations of X. Assume given a set R of equivalence relations on X, an element \(x_{0} \in X\) and a map \(\varphi : H \mapsto s_{H}\) from R to W. Denote by \(R_{0}\) the set of \(H \in R\) such that \(s_{\mathrm{H}}(x_{0}) \equiv x_{0} \mod \mathrm{H}'\) for all \(H' \neq H\) in R, and by \(S_{0}\) the set of \(s_{H}\) for H in \(R_{0}\) . We make the following assumptions:
\end{enumerate}

\begin{enumerate}
\def\labelenumi{(\roman{enumi})}
\item
  For any \(\mathbf{H} \in \Re\) , there are two equivalence classes modulo \(\mathbf{H}\) that are permuted by \(s_{\mathrm{H}}\) and \(s_{\mathrm{H}}^{2} = 1\) .
\item
  For all \(\mathbf{H} \in \Re\) and all \(w \in \mathbf{W}\) , the transform \(w(\mathbf{H})\) of \(\mathbf{H}\) by \(w\) is an equivalence relation belonging to \(\Re\) and \(s_{w(\mathbf{H})} = ws_{\mathbf{H}}w^{-1}\) .
\item
  For any \(w \neq 1\) in \(W\) , the set of \(H \in \Re\) such that \(w(x_0) \not\equiv x_0 \mod H\) is finite and meets \(\Re_0\) .
\end{enumerate}

\begin{enumerate}
\def\labelenumi{\alph{enumi})}
\item
  Prove that \((\mathbf{W},\mathbf{S}_0)\) is a Coxeter system (use Prop. 6 of no. 7).
\item
  Prove that the length \(l_{S_0}(w)\) is equal to the number of elements \(H \in \mathfrak{R}\) such that
\end{enumerate}

\[
w (x _ {0}) \not \equiv x _ {0} \bmod . H.
\]

\begin{enumerate}
\def\labelenumi{\alph{enumi})}
\setcounter{enumi}{2}
\tightlist
\item
  Let E be a finite set and X the set of total order relations on X. Denote by W the group of permutations of E, acting in the obvious way on E. Let i and j be distinct elements of E; say that two elements R and R' of X are equivalent mod. \(H_{ij}\) if either \(\mathrm{R}(i,j)\) and \(\mathrm{R}'(i,j)\) or \(\mathrm{R}(j,i)\) and \(\mathrm{R}'(j,i)\) ; and denote the transposition of i and j by \(s_{ij}\) . Let R be the set of equivalence relations of the form \(H_{ij}\) and \(\varphi\) the map from R to W defined by \(\varphi(\mathrm{H}_{ij}) = s_{ij}\) ; finally let \(x_{0}\) be an arbitrary element of X. Show that the objects thus defined satisfy assumptions (i) to (iii), and recover the results of Exerc. 4.
\end{enumerate}

\begin{enumerate}
\def\labelenumi{\arabic{enumi})}
\setcounter{enumi}{5}
\item
  Let E be a set with 6 elements and F the set of structures of the projective line over the field \(F_{5}\) on E. Denote by \(S_{E}\) the group of permutations of E; for any \(\sigma \in S_{E}\) denote by \(\tilde{\sigma}\) the permutation of F induced by \(\sigma\) by transport of structure. Show that there exists a bijection u from E to F, and that the map \(\sigma \mapsto u^{-1}\tilde{\sigma}u\) is an outer automorphism of \(S_{E}\) (if s is a transposition, \(u^{-1}\tilde{s}u\) has three orbits of two elements).
\item
  Construct a group W and two subsets S and S' of W such that (W, S) and (W, S') are isomorphic Coxeter systems, but such that there exists no inner automorphism of W transforming S to S' (use Exerc. 4 and Exerc. 6).
\item
  Construct a group W and two subsets S and S' of W such that (W, S) and (W, S') are non-isomorphic Coxeter systems, one of them being irreducible and the other not (for W take a dihedral group of order 12 generated by \(\{s, s'\}\) where s and \(s'\) are of order 2 and \(s \neq s'\) , and put \(S = \{s, s'\}\) and \(S' = \{(ss')^3, s', s'(ss')^2\}\) ).
\item
  Let (W, S) be a Coxeter system with matrix \((m(s, s'))\) , and let \(W^{+}\) be the subgroup of W consisting of the elements of even length. Let \(s_{0} \in S\) . Put \(g_{s} = ss_{0}\) . Show that the family \((g_{s})_{s \in S - \{s_{0}\}}\) and the relations \(g_{s}^{m(s, s_{0})} = 1\) for \(m(s, s_{0}) \neq \infty\) and \((g_{s} g_{s'}^{-1})^{m(s, s')} = 1\) for \(s, s' \in S - \{s_{0}\}\) and \(m(s, s') \neq \infty\) , form a presentation of \(W^{+}\) . (Let \(H^{+}\) be the group defined by the above presentation. Show that there exists an automorphism \(\alpha\) of \(H^{+}\) , with square the identity, that transforms \(g_{s}\) to \(g_{s}^{-1}\) for all \(s \in S - \{s_{0}\}\) . If \(H_{\alpha}\) is the semi-direct product of \(\{1, -1\}\) and \(H^{+}\) , relative to \(\alpha\) , define mutually inverse homomorphisms \(H_{\alpha} \to W\) and \(W \to H_{\alpha}\) .) Show that if the elements of S are conjugate (cf.~Prop. 3), the group \(W^{+}\) is the commutator subgroup of W (remark that the elements \(g_{s}\) are then commutators).
\item
  Let \(\mathfrak{U}_n\) be the alternating group consisting of the permutations \(w \in \mathfrak{S}_n\) whose signature is equal to \(+1\) . Show that \(\mathfrak{U}_n\) is the commutator subgroup of \(\mathfrak{S}_n\) (use Exerc. 4 and 9). For any integer \(i\) with \(1 \leqslant i \leqslant n - 2\) , put \(u_i = s_i s_{i+1}\) (in the notation of Exerc. 4). Show that the family \((u_i)\) and the relations \(u_1^3 = 1\) , \(u_i^2 = 1\) for \(i \geqslant 2\) , \((u_i u_{i+1})^3 = 1\) for \(1 \leqslant i \leqslant n - 3\) , and \(u_i u_j = u_j u_i\) for \(1 \leqslant i \leqslant n - 4\) and \(i + 2 \leqslant j \leqslant n - 2\) , form a presentation of the group \(\mathfrak{U}_n\) (use Exerc. 9).
\end{enumerate}

*11) Let \((\mathbf{W},\mathbf{S})\) be a Coxeter system. Let \(\Gamma_{\infty}\) be the graph whose set of vertices is \(\mathbf{S}\) , two vertices \(s,s'\) being joined by an edge if and only if \(m(s,s')\neq \infty\) . Let \(\mathrm{S}_{\alpha}\) be the connected components of \(\Gamma_{\infty}\) . Show that \(\mathbf{W}\) can be identified with the free product of the \(\mathrm{W}_{\mathrm{S}_{\alpha}}\) . In particular, every \(w\in \mathbf{W}\) can be written uniquely as a product \(w_{1}\ldots w_{h}\) with \(w_{i}\neq 1\) , \(w_{i}\) belonging to \(\mathrm{W}_{\mathrm{S}_{\alpha_i}}\) , and \(\alpha_{i}\neq \alpha_{i + 1}\) for \(1\leqslant i\leqslant h - 1\) ; show that the length of \(w\) is the sum of the lengths of the \(w_i\).*

\begin{enumerate}
\def\labelenumi{\arabic{enumi})}
\setcounter{enumi}{11}
\tightlist
\item
  Let \((\mathbf{W},\mathbf{S})\) be a Coxeter system such that \(m(s,s')\) is even if \(s\neq s'\) and let \(\mathbf{X}\) be a subset of \(\mathbf{S}\) . Show that there exists a unique homomorphism \(f_{\mathbf{X}}\) from \(\mathbf{W}\) to \(\mathbf{W}_{\mathbf{X}}\) such that \(f_{\mathbf{X}}(s) = s\) for \(s\in \mathbf{X}\) and \(f_{\mathbf{X}}(s) = 1\) for \(s\in \mathbf{S} - \mathbf{X}\) . Deduce that \(\mathbf{W}\) is the semi-direct product of \(\mathbf{W}_{\mathbf{X}}\) and the kernel of \(f_{\mathbf{X}}\) . Show that if \(\mathbf{X}\subset \mathbf{Y}\subset \mathbf{S}\) , there exists a unique homomorphism \(f_{\mathbf{X},\mathbf{Y}}\) from \(\mathbf{W}_{\mathbf{Y}}\) to \(\mathbf{W}_{\mathbf{X}}\) such that \(f_{\mathbf{X}} = F_{\mathbf{X},\mathbf{Y}}\circ f_{\mathbf{Y}}\) and that \(\mathbf{W}\) can be identified with a subgroup of the projective system thus obtained from the \(W_{X}\) as X runs through the filtered set of finite subsets of S.
\end{enumerate}

¶ 13) Let \((W, S)\) be a Coxeter system. For \(s, s' \in S\) , define the sequence \(\mathbf{a}(s, s')\) by means of the following rules:

\begin{enumerate}
\def\labelenumi{(\roman{enumi})}
\item
  if \(ss'\) is of infinite order, \(\mathbf{a}(s, s')\) is the empty sequence;
\item
  if \(ss'\) is of finite order m the sequence \(\mathbf{a}(s, s')\) is of length m and its even (resp. odd) numbered terms are equal to \(s'\) (resp. s).
\end{enumerate}

Denote by \(a(s, s')\) the product of the sequence \(\mathbf{a}(s, s')\) .

\begin{enumerate}
\def\labelenumi{\alph{enumi})}
\item
  Show that the generating set S and the relations \(s^{2}=1\) and \(\mathbf{a}(s,s')=\mathbf{a}(s',s)\) form a presentation of the group W.
\item
  Let \(q\) be an integer \(\geqslant 1\) . Let \(\Sigma_q\) be the set of sequences of \(q\) elements of \(S\) and let \(R_q\) be the smallest equivalence relation on \(\Sigma_q\) for which sequences of the form \((A, \mathbf{a}(s, s'), B)\) and \((A, \mathbf{a}(s', s), B)\) (where \(s, s' \in S\) and \(A\) and \(B\) are sequences of elements of \(S\) ) are equivalent. Let \(\Sigma_q^r\) be the set of sequences \(\mathbf{s} = (s_1, \ldots, s_q)\) such that \(w(\mathbf{s}) = s_1 \ldots s_q\) is of length \(q\) . Show that sequences \(\mathbf{s}, \mathbf{s}' \in \Sigma_q^r\) are equivalent modulo \(R_q\) if and only if \(w(\mathbf{s}) = w(\mathbf{s}')\) (argue by induction on \(q\) and apply Prop. 5 of no. 5).
\item
  Show that a sequence \(\mathbf{s} \in \Sigma_q\) does not belong to \(\Sigma_q^r\) if and only if \(\mathbf{s}\) is equivalent modulo \(R_q\) to a sequence in which two consecutive terms are equal. (Argue by induction on \(q\) and reduce to the case of a sequence \((s_1, \ldots, s_q)\) which does not belong to \(\Sigma_q^r\) but which is such that \((s_1, \ldots, s_{q-1})\) and \((s_2, \ldots, s_q)\) belong to \(\Sigma_{q-1}^r\) ; use Exerc. 1 to show that \(s_1 \ldots s_{q-1} = s_2 \ldots s_q\) and apply b).)
\end{enumerate}

*14) Let \((\mathbf{W},\mathbf{S})\) be a Coxeter system and let \((\Gamma ,f)\) be its Coxeter graph. Let \(k\) be an integer \(\geqslant 3\) . If \(a\) is an edge of \(\Gamma\) , put \(f_{k}(a) = f(a)\) if \(f(a)\neq \infty\) and \(f_{k}(a) = k\) if \(f(a) = \infty\) . Let \((\mathrm{W}_k,\mathrm{S})\) be a Coxeter system whose Coxeter graph is equal to \((\Gamma ,f_k)\) (Chap. V, §4, no. 3, Cor. of Prop. 4). Show that there exists a unique homomorphism \(\varphi_{k}\) from W to \(\mathbf{W}_k\) inducing the identity on S. Show that if \(k\) divides \(k^{\prime}\) , there exists a unique homomorphism \(\varphi_{k,k^{\prime}}\) from \(\mathbf{W}_{k^{\prime}}\) to \(\mathbf{W}_k\) such that \(\varphi_{k} = \varphi_{k,k^{\prime}}\circ \varphi_{k^{\prime}}\) . Show that the homomorphism \((\varphi_{k})\) from W to the projective limit of the \(\mathbf{W}_k\) is injective (use Exerc. 13) \(c)\) , but in general not surjective (for example, in the case of the infinite dihedral group).*

\begin{enumerate}
\def\labelenumi{\arabic{enumi})}
\setcounter{enumi}{14}
\tightlist
\item
  Let A be a set and C a subset of \(\mathfrak{P}(A)\) . The elements of C are called the chambers of A and a subset F of a chamber C is called a facet, the codimension of F in C being the cardinal of C-F. A facet F is said to be a panel of C if F is of codimension 1 in C. Two chambers C and C' are said to be adjoining if they have a common panel F: then either \(C = C'\) or \(F = C \cap C'\) . A gallery of length n is a sequence \(\Gamma = (C_0, C_1, \ldots, C_n)\) of \(n + 1\) chambers such that \(C_i\) and \(C_{i+1}\) are adjoining for \(0 \leqslant i \leqslant n - 1\) . Then \(C_0\) and \(C_n\) are called the ends of \(\Gamma\) . The gallery \(\Gamma\) is called injective if \(C_i \neq C_{i+1}\) for \(0 \leqslant i \leqslant n - 1\) and is called minimal if there is no gallery with the same ends and length \textless{} n.
\end{enumerate}

The set A (together with C) is a building if every element of A belongs to at least one chamber and if any two chambers are the ends of a gallery. The distance between two chambers C and \(C'\) is the length \(d(\mathrm{C},\mathrm{C}')\) of a minimal gallery with ends C and \(C'\) .

A sub-building of a building A is a subset D of A such that D together with \(\mathfrak{C}\cap\mathfrak{P}(D)\) is a building.

\begin{enumerate}
\def\labelenumi{\alph{enumi})}
\item
  Show that if A is a building, a facet has the same codimension in every chamber containing it; this allows us to speak of the codimension of a facet and of a panel of A without reference to a particular chamber. A morphism of a building B to A is a map f from B to A such that the restriction of f to every chamber C of B is a bijection from C to a chamber \(f(\mathrm{C})\) of A. Show that the image under f of a facet of B is a facet of the same codimension.
\item
  A building is called an apartment if every panel is contained in exactly two chambers. Show that if A is an apartment, every automorphism of A (i.e.~every permutation of A preserving C) that leaves fixed all the points of a chamber is the identity. More generally, let \(\varphi\) be an endomorphism of A and let C be a chamber of A such that \(\varphi(a)=a\) for all \(a\in C\) . Let \((\mathrm{C},\mathrm{C}_{1},\ldots,\mathrm{C}_{n})\) be a gallery of A. Show that, either the gallery \((\mathrm{C},\varphi(\mathrm{C}_{1}),\ldots,\varphi(\mathrm{C}_{n}))\) is not injective, or \(\varphi(a)=a\) for all a belonging to the union of the \(C_{i}\) .
\end{enumerate}

\begin{enumerate}
\def\labelenumi{\arabic{enumi})}
\setcounter{enumi}{15}
\tightlist
\item
  Let \((W,S)\) be a Coxeter system. For \(s \in S\) , denote by \(W^{(s)}\) the subgroup \(W_{S-\{s\}}\) of W, by A the set of subsets of W of the form \(wW^{(s)}\) for \(w \in W\) and \(s \in S\) , and by C the set of subsets of A of the form \(C_w = \{wW^{(s)} \mid s \in S\}\) for \(w \in W\) , which we call the chambers of A (Exerc. 15).
\end{enumerate}

\begin{enumerate}
\def\labelenumi{\alph{enumi})}
\item
  Show that the map \(w \mapsto \mathbf{C}_w\) is a bijection from W onto \(\mathfrak{C}\) .
\item
  Show that two distinct chambers \(C_{w}\) and \(C_{w'}\) are adjoining if and only if there exists \(s \in S\) such that \(w' = ws\) . Deduce that A (together with C) is an apartment (Exerc. 15), which we call the apartment of (W, S). Show that the length of a minimal gallery with ends \(C_{w}\) and \(C_{w'}\) is equal to \(l_{S}(w^{-1}w')\) .
\item
  Let \(\mathfrak{F}\) be the set of facets of A and let \(F \in \mathfrak{F}\) . Show that there exist a unique subset X of S and an element \(w \in W\) such that \(wW_X = \bigcap_{i \in F} i\) . Then F is said to be of type X. Show that the codimension of F is equal to the cardinal of X. Show that the map \(j: F \mapsto \bigcap_{i \in F} i\) is a strictly decreasing bijection (with respect to inclusion) from \(\mathfrak{F}\) to the set of subsets of W of the form \(wW_X\) for \(w \in W\) and \(X \subset S\) . Show that every facet of type X contains a unique facet of type Y for every Y such that \(X \subset Y \subset S\) .
\item
  Let W act on A be left translations and put \(C = C_{e}\) . Show that W acts simply-transitively \(^{8}\) on C. Let \(C_{1}, \ldots, C_{n}\) be chambers of A. Show that the following conditions are equivalent:
\end{enumerate}

\begin{enumerate}
\def\labelenumi{(\roman{enumi})}
\item
  the sequence \(\Gamma = (\mathrm{C}_0 = \mathrm{C},\mathrm{C}_1,\dots ,\mathrm{C}_n)\) is an injective gallery;
\item
  there exists a sequence \(\mathbf{s} = (s_1, \ldots, s_n)\) of elements of S such that \(\mathrm{C}_j = t_j(\mathrm{C}_{j-1})\) for \(1 \leqslant j \leqslant n\) , where \(t_j\) is the element of \(\Phi(\mathbf{s})\) defined in no. 4, formula (11).
\end{enumerate}

Show that if these conditions are satisfied, the sequence-s is unique and \(C_{n} = s_{1} \ldots s_{n}(C)\) . Show that the gallery \(\Gamma\) is minimal if and only if the sequence \(\mathbf{s}(\Gamma) = \mathbf{s}\) is a reduced decomposition of \(w = s_{1} \ldots s_{n}\) .

\begin{enumerate}
\def\labelenumi{\alph{enumi})}
\setcounter{enumi}{4}
\tightlist
\item
  Let T be the union of the conjugates of S. For \(t \in T\) , the set \(L_{t}\) of points of A invariant under t is called the wall defined by t. Show that \(L_{t}\) is a union of panels and that a panel F is contained in \(L_{t}\) if and only if \(j(\mathrm{F})\) is of the form \(wW_{\{s\}}\) with \(t = wsw^{-1}\) . Deduce that, for any panel F, there is a unique element \(t = t(\mathrm{F}) \in \mathrm{T}\) such that \(F \subset L_{t}: L_{t}\) is called the support of F.
\end{enumerate}

Show that if \(w(\mathrm{L}_t) = \mathrm{L}_t\) (for \(w \in \mathrm{W}\) ), then \(w = 1\) or \(w = t\) .

\(f)\) Let \(w', w'' \in W\) . Put \(C' = w'(C), C'' = w''(C)\) and let \(\Gamma = (C_0 = C', C_1, \ldots, C_n = C'')\) be an injective gallery with ends \(C'\) and \(C''\) . Let \(t_j\) be the element of \(T\) defining the support wall of the panel \(C_j \cap C_{j-1}\) (for \(1 \leqslant j \leqslant n\) ). Show that the sequence \(\Psi(T) = (w'^{-1} t_j w')_{1 \leqslant j \leqslant n}\) coincides with the sequence \(\Phi(s(w'^{-1}(\Gamma))\) . For \(t \in T\) , let \(n(\Gamma, t)\) be the number of times \(w'^{-1} tw'\) appears in \(\Psi(\Gamma)\) . Deduce from Lemma 1 of no. 4 that the number \((-1)^{n(\Gamma, t)}\) depends only on \(C'\) and \(C''\) and not on \(\Gamma\) : we denote it by \(\eta(C', C'', t)\) . Show that the relation \(\eta(C', C'', t) = 1\) is an equivalence relation between \(C'\) and \(C''\) that has two equivalence classes that are permuted by \(t\) . Denote by \(\mathfrak{C}^+(t)\) the equivalence class containing \(C\) and by \(\mathfrak{C}^-(t)\) the other.

Show that, for \(w \in \mathbf{W}\) and \(s \in S\) , the chamber \(w(\mathbf{C})\) belongs to \(\mathfrak{C}^+(s)\) if and only if \(l(sw) > l(w)\) .

\(g)\) Let \(\mathrm{A}^{+}(t)\) (resp. \(\mathrm{A}^{-}(t)\) ) be the union of the chambers belonging to \(\mathfrak{C}^{+}(t)\) (resp. \(\mathfrak{C}^{-}(t)\) ) (for \(t \in \mathrm{T}\) ). Show that \(\mathrm{A}^{+}(t) \cap \mathrm{A}^{-}(t) = \mathrm{L}_{t}\) . (To show that \(\mathrm{A}^{+}(t) \cap \mathrm{A}^{-}(t) \subset \mathrm{L}_{t}\) , reduce to the case \(t \in \mathrm{S}\) . If \(a \in \mathrm{A}^{+}(t) \cap \mathrm{A}^{-}(t)\) , put \(a = w\mathrm{W}^{(s)}\) with \(s \in \mathrm{S}\) and \(w\) being \((\varnothing, \mathrm{S} - \{s\})\) -reduced (Exerc. 3). If \(w(\mathrm{C}) \in \mathfrak{C}^{-}(t)\) , then \(l(tw) < l(w)\) and \(w = ts_{1}, \ldots, s_{q}\) with \(s_{j} \in \mathrm{S}\) and \(l(w) = q + 1\) . Since \(a \in \mathrm{A}^{-}(t)\) , there exists \(x \in \mathrm{W}^{(s)}\) such that \(wx(\mathrm{C}) \in \mathfrak{C}^{+}(t)\) . Then

\[
l (t w x) = 1 + l (w x) = 1 + l (w) + l (x),
\]

but since \(twx = s_1 \ldots s_q x\) this is a contradiction. Thus, \(w(\mathbf{C}) \in \mathfrak{C}^+(t)\) and \(l(tw) = 1 + l(w)\) . If \(x \in \mathrm{W}^{(s)}\) is such that \(l(twx) < l(wx)\) , deduce from Exerc. 1 that \(twx = wx'\) with \(x' \in \mathrm{W}^{(s)}\) , and hence that \(ta = a\) and \(a \in \mathrm{L}_t\) .

The subsets \(\mathbf{A}^{+}(t)\) and \(\mathbf{A}^{-}(t)\) are called the halves of A determined by the wall \(L_{t}\) . Two points of A are said to be on the same side (resp. strictly on opposite sides) of \(L_{t}\) if they belong (resp. do not both belong) to one of the two halves. Any facet is contained in one of the halves determined by \(L_{t}\) . If two facets are contained in different halves, they are said to be on opposite sides of \(L_{t}\) , or to be separated by \(L_{t}\) .

\(h\) ) Let \(w \in \mathbf{W}\) . Show that \(l(w)\) is equal to the number of walls separating \(\mathbf{C}\) and \(w(\mathbf{C})\) .

\begin{enumerate}
\def\labelenumi{\roman{enumi})}
\tightlist
\item
  Show that the map \(\varphi\) that takes the half \(A^{+}(t)\) (resp. \(A^{-}(t)\) ) to \((1,t)\) (resp. \((-1,t)\) ) is a bijection from the set \(\mathfrak{M}\) of halves of \(A\) onto the set \(R = \{1, - 1\} \times T\) (cf.~no. 4). With the notation of Lemma 1 of no. 4, show that \(\varphi(w(M)) = U_w(\varphi(M))\) for all \(w\in W\) and \(M\in \mathfrak{M}\) .
\end{enumerate}

\begin{enumerate}
\def\labelenumi{\arabic{enumi})}
\setcounter{enumi}{16}
\item
  We retain the notation of Exerc. 16 and assume that W is finite. Let H be the set of walls of A. To any \(H \in H\) , we associate the half \(H^{+}\) of A determined by H containing C. Show that the elements of H can be numbered in such a way that the map \(j \mapsto \bigcap_{i \leqslant j} H_{i}^{+}\) is strictly decreasing on the interval {[}1, \(\operatorname{Card}(\mathfrak{H})\) {]}. (Consider the family \(\mathfrak{F}\) of intersections of the sets \(\mathrm{H}^{+}\) ordered by inclusion and consider a strictly decreasing sequence \((\mathrm{F}_{0},\ldots,\mathrm{F}_{q})\) of elements of \(\mathfrak{F}\) of maximal length. For any \(\mathrm{H}\in\mathfrak{H}\) , there exists an \(i\) such that \(\mathrm{H}^{+}\supset\mathrm{F}_{j}\) for \(j\geqslant i\) and \(\mathrm{H}^{+}\not\supset\mathrm{F}_{i-1}\) : show that \(\mathrm{F}_{i}=\mathrm{H}^{+}\cap\mathrm{F}_{i-1}\) .)
\item
  Let A be an apartment (Exerc. 15). A folding of A is an endomorphism \(\pi\) of A such that \(\pi^{2} = \pi\) and such that every chamber of A is the image under \(\pi\) of 0 or 2 chambers.
\end{enumerate}

\begin{enumerate}
\def\labelenumi{\alph{enumi})}
\item
  Let \(\pi\) be a folding of A and C a chamber of A such that \(\pi(C) = C\) . For any neighbouring chamber \(C'\) of C, either \(\pi(C') = C'\) or \(\pi(C') = C\) ; if \(a \in C\) we have \(\pi(a) = a\) . Show that, if \((C_0 = C, C_1, \ldots, C_n)\) is a gallery, either \(\pi(C_i) = C_i\) for all \(i\) or \((C_0, \pi(C_1), \ldots, \pi(C_n))\) is not minimal (and in fact has two consecutive equal chambers). Deduce that any minimal gallery whose extremities are invariant under \(\pi\) is invariant under \(\pi\) . If \((C = C_0, C_1, \ldots, C_n)\) is a minimal gallery and if \(\pi(C_n) \neq C_n\) , there exists an index \(i\) with \(0 \leqslant i < n\) such that \(\pi(C_j) = C_j\) for \(0 \leqslant j \leqslant i\) and \(\pi(C_j) \neq C_j\) for \(i < j \leqslant n\) .
\item
  Let \(C_{1}\) and \(C_{2}\) be two distinct neighbouring chambers and \(\pi, \pi'\) two foldings of A. Assume that \(\pi(C_{0}) = C_{1}\) and \(\pi'(C_{1}) = C_{0}\) . Let C be a chamber, and consider the following three conditions:
\end{enumerate}

\[
\pi (\mathrm{C}) = \mathrm{C};\tag{1}
\]

\[
d (\mathrm{C}, \mathrm{C} _ {1}) <   d (\mathrm{C}, \mathrm{C} _ {2});\tag{2}
\]

\[
\pi^ {\prime} (\mathrm{C}) \neq \mathrm{C}.\tag{3}
\]

Show that (1) \(\Longrightarrow\) (2) \(\Longrightarrow\) (3) and deduce that the three conditions are equivalent. Show that \(\pi\) (resp. \(\pi'\) ) is the unique folding of A taking \(C_2\) (resp.

\(C_{1}\) to \(C_{1}\) (resp. \(C_{2}\) ) (assume that (2) is satisfied and let \((C_{1}, C_{2}', \ldots, C_{n}' = C)\) be a minimal gallery: show that \(\pi'(C_{j+1}')\) is the unique chamber distinct from \(\pi'(C_{j}')\) and containing the panel \(\pi'(C_{j}' \cap C_{j+1}')\) . Show that \(\pi(\mathfrak{C})\) and \(\pi'(\mathfrak{T})\) form a partition of the set T of chambers of A and that \(\pi(a) = \pi'(a) = a\) for all \(a \in \pi(A) \cap \pi'(A)\) . Show that the map from A to itself that coincides with \(\pi'\) on \(\pi(A)\) and with \(\pi\) on \(\pi'(A)\) is an involutive automorphism of A. It is called the reflection with respect to the panel \(C_{1} \cap C_{2}\) . Show that it is the only non-trivial automorphism of A leaving fixed all the points of \(C_{1} \cap C_{2}\) (use Exerc. 15 b)).

\begin{enumerate}
\def\labelenumi{\alph{enumi})}
\setcounter{enumi}{2}
\tightlist
\item
  Suppose that A is the apartment associated to a Coxeter system (W,S) and retain the notation of Exerc. 16. Let \(C_{1}\) and \(C_{2}\) be two neighbouring chambers and let t be the element of T such that the wall \(L_{t}\) is the support of \(C_{1} \cap C_{2}\) . Let \(M_{j}\) be the half of A determined by \(L_{t}\) and containing \(C_{j}\) (for j = 1,2). Show that the map \(\pi\) defined by \(\pi(a) = a\) if \(a \in M_{1}\) and \(\pi(a) = t(a)\) if \(a \in M_{2}\) is a folding of A such that \(\pi(C_{2}) = C_{1}\) and that the reflection with respect to the panel \(C_{1} \cap C_{2}\) is the map \(a \mapsto t(a)\) .
\end{enumerate}

\begin{enumerate}
\def\labelenumi{\arabic{enumi})}
\setcounter{enumi}{18}
\tightlist
\item
  Let A be an apartment. Assume that, for any two distinct neighbouring chambers \(C_{1}\) and \(C_{2}\) , there is a folding (Exerc. 18) of A taking \(C_{1}\) to \(C_{2}\) . Let C be a chamber of A and \((\mathrm{C}_{i})_{i\in I}\) the family of chambers neighbouring C and distinct from C. Denote by \(s_{i}\) the reflection with respect to the panel \(C\cap C_{i}\) (Exerc. 18 b)). Put \(S=\{s_{i}\mid i\in I\}\) and denote by W the group of automorphisms of A generated by the \(s_{i}\) .
\end{enumerate}

\begin{enumerate}
\def\labelenumi{\alph{enumi})}
\item
  Show that, for any chamber \(C'\) , there exists \(w \in W\) such that \(C' = w(C)\) (argue by induction on \(d(C, C')\) ).
\item
  Show that (W,S) is a Coxeter system (for \(i\in I\) ,put
\end{enumerate}

\[
\mathrm{P} _ {s _ {i}} = \left\{w \in \mathrm{W} \mid w (\mathrm{C}) \subset \pi_ {i} (\mathrm{A}) \right\},
\]

where \(\pi_{i}\) is the folding taking \(C_{i}\) to C, and show that the assumptions of Prop. 6 are satisfied: to prove condition (C'), remark that if \(w \in P_{s_{i}}\) and \(ws_{j} \notin P_{s_{i}}\) , then

\[
w (\mathrm{C}) \cap w s _ {j} (\mathrm{C}) \subset \pi_ {i} (\mathrm{A}) \cap s _ {i} \pi_ {i} (\mathrm{A}).
\]

Since \(w(\mathbf{C})\) and \(ws_{j}(\mathbf{C})\) are neighbouring, it follows that \(s_i = ws_jw^{-1}\) (Exerc. 18 b)).

\begin{enumerate}
\def\labelenumi{\alph{enumi})}
\setcounter{enumi}{2}
\item
  Let F be a facet of the chamber C. Show that the stabilizer \(W_{F}\) of F in W is generated by the \(s_{i} \in S\) such that \(F \subset C \cap C_{i}\) (let \(w \in W_{F}\) with \(l_{S}(w) > 1\) and let \(i \in I\) be such that \(w = s_{i}w'\) with \(l(w') = l(w) - 1\) ; by Prop. 6, \(w' \in P_{s_{i}}\) , hence \(w(\mathrm{C}) \subset s_{i}\pi_{i}(\mathrm{A})\) , \(F \subset \pi_{i}(\mathrm{A}) \cap s_{i}\pi_{i}(\mathrm{A})\) and \(s_{i} \in W_{F}\) ). In particular, \(w(\mathrm{C}) = \mathrm{C}\) if and only if w = 1.
\item
  Show that the map \(a \mapsto W_{\{a\}}\) is an isomorphism from the apartment A to the apartment associated to (W, S) (Exerc. 16), compatible with the action of W.
\end{enumerate}

\begin{enumerate}
\def\labelenumi{\arabic{enumi})}
\setcounter{enumi}{19}
\tightlist
\item
  Let A be a building and S a set. Then A is said to be numbered by S if one is given a map f from A to S such that, for any chamber C of A, the restriction of f to C is a bijection from C to S. If F is a facet of A, \(f(F)\) is called the type of F. Let A be a numbered building. An endomorphism \(\varphi\) of A is called allowed if a and \(\varphi(a)\) are of the same type for all \(a \in A\) .
\end{enumerate}

\begin{enumerate}
\def\labelenumi{\alph{enumi})}
\item
  Let \(\varphi\) be an endomorphism of A. Show that, if there exists a chamber C of A such that a and \(\varphi(a)\) are of the same type for all \(a \in C\) , then \(\varphi\) is allowed. Show that if A is an apartment and \(\pi\) is a folding of A (Exerc. 18), then \(\pi\) is an allowed endomorphism.
\item
  A subset D of A is called convex if, for all \(a \in A - D\) , there exists an allowed endomorphism \(\varphi\) of A such that \(\varphi(x) = x\) for all \(x \in D\) and \(\varphi(a) \neq a\) . Show that any intersection of convex sets is convex and that, for any subset D of A, there exists a smallest convex subset containing D: this is called the convex hull of D and denoted by \(\Gamma(D)\) .
\end{enumerate}

\begin{enumerate}
\def\labelenumi{\arabic{enumi})}
\setcounter{enumi}{20}
\tightlist
\item
  Let \((W,S)\) be a Coxeter system and A the associated apartment (cf.~Exerc. 16, of which we retain the notation).
\end{enumerate}

\begin{enumerate}
\def\labelenumi{\alph{enumi})}
\item
  Show that there exists a unique numbering of A (called the canonical numbering) for which the type of a facet F is that defined in Exerc. 16 c). We shall always consider A to be equipped with this numbering.
\item
  Show that the allowed automorphisms of \(\mathbf{A}\) are the operations of \(\mathbf{W}\) .
\item
  Let D be a subset of A containing at least one chamber. Show that the following conditions are equivalent:
\end{enumerate}

\begin{enumerate}
\def\labelenumi{(\roman{enumi})}
\item
  D is the intersection of the halves of A (Exerc. 16 g)) that contain D;
\item
  D is convex;
\item
  whenever two facets \(F_{1}\) and \(F_{2}\) are contained in D, the convex hull of \(F_{1} \cup F_{2}\) is contained in D;
\item
  whenever a chamber \(\mathbf{C}_1\) and a facet \(\mathbf{F}\) are contained in \(\mathrm{D}\) and \((C_1, \ldots, C_n)\) is a gallery of minimal length such that \(\mathbf{F} \subset \mathbf{C}_n\) , we have \(C_i \subset D\) for \(1 \leqslant i \leqslant n\) .
\end{enumerate}

(To show that (iii) \(\Longrightarrow\) (iv), use Exerc. 15 \(b\) ). To show that (iv) \(\Longrightarrow\) (i), argue by contradiction. Let \(D'\) be the intersection of the halves of A containing D; let \(a \in D' - D\) , let \(C_0\) be a chamber contained in D and let \((C_0, C_1, \ldots, C_n)\) be a gallery of minimal length such that \(a \in C_n\) . Then \(C_j \subset D'\) for all \(j\) . Show that there exists an integer \(j\) with \(0 \leqslant j < n\) such that \(C_j \subset D\) and \(C_{j+1} \not\subset D\) . Let M (resp. \(M'\) ) be the half of A determined by the support wall of the panel \(C_j \cap C_{j+1}\) and containing \(C_j\) (resp. \(C_{j+1}\) ). Show that D \(\not\subset M\) . Let \(b \in D \cap (A - M)\) and let \(\Gamma = (C_j, C_1', \ldots, C_p')\) be a gallery of smallest possible length such that \(b \in C_p'\) . Then \(C_k' \subset D\) for \(1 \leqslant k \leqslant p\) and \(C_p' \subset M'\) . Let \(\pi\) be the folding of A with image \(M'\) (Exerc. 18 \(c\) ): then \(\pi(C_j) = C_{j+1}\) and the gallery \(\pi(\Gamma)\) is not injective (Exerc. 18 \(a\) ). If \(\Gamma' = (C_{j+1}, C_2', \ldots, C_{p-2}', C_p')\)

is the gallery obtained from \(\pi(\Gamma)\) by suppressing one of the two equal consecutive chambers, the gallery \((\mathrm{C}_{j}, \mathrm{C}_{j+1}, \mathrm{C}_{2}', \ldots, \mathrm{C}_{p-2}', \mathrm{C}_{p}')\) is minimal by the definition of \(\Gamma\) . Deduce from (iv) that \(\mathrm{C}_{j+1} \subset \mathrm{D}\) . Contradiction.)

\begin{enumerate}
\def\labelenumi{\arabic{enumi})}
\setcounter{enumi}{21}
\tightlist
\item
  We retain the notation of Exerc. 16 and 21.
\end{enumerate}

\begin{enumerate}
\def\labelenumi{\alph{enumi})}
\item
  Let \(t \in \mathbf{T}\) and \(w \in \mathbf{W}\) . Show that the chambers \(\mathbf{C}\) and \(w(\mathbf{C})\) are separated by the wall \(\mathbf{L}_t\) if and only if \(l(tw) < l(w)\) (use the folding determined by the half \(\mathbf{A}^+(t)\) ).
\item
  Let \(w_0 \in W\) . Show that the following conditions are equivalent:
\end{enumerate}

\[
l (w w _ {0}) = l (w _ {0}) - l (w) \quad \text {   for   all   } w \in \mathrm{W};\tag{i}
\]

\[
l (t w _ {0}) <   l (w _ {0}) \quad \text {   for   all   } t \in \mathrm{T};\tag{ii}
\]

whenever \(t \in \mathrm{T}\) , the chambers \(\mathrm{C}\) and \(w_0(\mathrm{C})\) are separated by the wall \(\mathbf{L}_t\) (iii)

(Use Exerc. 16 h) to show that (iii) \(\Longrightarrow\) (i).)

Show that such an element \(w_{0}\) is unique and exists if and only if W is finite. It is then the element of greatest length of W and is characterized by

\[
l (s w _ {0}) <   l (w _ {0}) \quad \text {   for   all   } s \in S.\tag{iv}
\]

Moreover, \(w_0^2 = 1\) , \(w_0Sw_0 = S\) and \(l(w_0) = \mathrm{Card}(T)\) .

\begin{enumerate}
\def\labelenumi{(\alph{enumi})}
\setcounter{enumi}{2}
\tightlist
\item
  Assume that W is finite. Show that, for any chamber \(C_{0}\) , there exists a unique chamber \(-C_{0}\) such that \(C_{0} \cup (-C_{0})\) is not contained in any half of A. Show that there exists a unique involutive automorphism \(\varphi\) of A (not necessarily allowed such that \(\varphi(C_{0}) = -C_{0}\) for any chamber \(C_{0}\) and that \(\varphi(L) = L\) for any wall L of A. We put \(\varphi(a) = -a\) for \(a \in A\) . If F is a facet, the facet \(-F = \varphi(F)\) will be said to be opposed to F.
\end{enumerate}

\begin{enumerate}
\def\labelenumi{\alph{enumi})}
\setcounter{enumi}{3}
\tightlist
\item
  Let \(C_{0}\) be a chamber of A and F a partition of \(C_{0}\) . Show that the convex hull of \(C_{0} \cup (-F)\) is the half of A determined by the wall L that is the support of F and that contains \(C_{0}\) .
\end{enumerate}

\begin{enumerate}
\def\labelenumi{\arabic{enumi})}
\setcounter{enumi}{22}
\item
  We retain the notation of Exerc. 16. Let \(\operatorname{Aut}(A)\) be the group of automorphisms of the apartment \(A\) . Show that if \(\varphi \in \operatorname{Aut}(A)\) , then \(\varphi\) permutes the walls of \(A\) , and \(\varphi t\varphi^{-1} \in T\) for all \(t \in T\) (use Exerc. 18). Deduce that \(W\) (identified with a subgroup of \(\operatorname{Aut}(A)\) ) is normal and that \(\operatorname{Aut}(A)\) is the semi-direct product of the subgroup \(E\) of automorphisms preserving the chamber \(C\) by \(W\) . Show that the action of \(\operatorname{Aut}(A)\) on \(W\) defines an isomorphism from \(E\) to the group of automorphisms of the Coxeter system \((W, S)\) , or to that of the Coxeter graph of \((W, S)\) (cf.~also Exerc. 19).
\item
  A structured building is a building I equipped with a set U of sub-buildings satisfying the following conditions:
\end{enumerate}

(SB 1) The sub-buildings \(A \in \mathfrak{U}\) are apartments.

(SB 2) Any two chambers of I are contained in at least one element of U. (SB 3) Whenever \(A_{1}, A_{2} \in U\) are such that \(A_{1} \cap A_{2}\) contains a chamber, there exists an isomorphism from \(A_{1}\) to \(A_{2}\) leaving fixed the points of \(A_{1} \cap A_{2}\) .

Let \((\mathbf{I},\mathfrak{U})\) be a structured building. The elements of \(\mathfrak{U}\) are called the apartments of \((\mathbf{I},\mathfrak{U})\) , or simply those of \(\mathbf{I}\) .

\begin{enumerate}
\def\labelenumi{\alph{enumi})}
\item
  Show that any two apartments of I are isomorphic. Let C be a chamber of I and A an apartment of I containing C. Show that there exists a unique endomorphism \(\rho\) of I (called the retraction with centre C of I onto A) such that \(\rho(a)=a\) for all \(a\in A\) and that, for any apartment \(A'\) containing C, the restriction of \(\rho\) to \(A'\) is an isomorphism from \(A'\) to A (remark that, by (SB 2) and Exerc. 15 b), for any apartment \(A'\) containing C, there exists a unique isomorphism \(\rho_{A'}\) of \(A'\) leaving all the points of C fixed). Show that \(\rho^{2}=\rho\) and that \(\rho^{-1}(C)=C\) .
\item
  Let A be an apartment of I, \(C_{0}\) a chamber and F a facet contained in A. Let \((C_{0}, C_{1}, \ldots, C_{n})\) be a gallery of smallest possible length such that \(F \subset C_{n}\) . Show that \(C_{i} \subset A\) for \(1 \leqslant i \leqslant n\) (argue by contradiction: if \(C_{i} \subset A\) and \(C_{i+1} \notin A\) , consider the retraction of I onto A with centre the chamber of A distinct from \(C_{i}\) and containing \(C_{i} \cap C_{i+1}\) ).
\item
  Let A be an apartment of I, C a chamber of A, F a panel of C, C' a chamber of I and \(\Gamma = (C_0, \ldots, C_n = C')\) a gallery of smallest possible length such that \(F \subset C_0\) . Show that the retraction \(\rho\) of I onto A with centre C transforms \(\Gamma\) into a gallery
\end{enumerate}

\[
\left(\mathrm{C} _ {0} ^ {\prime}, \dots , \mathrm{C} _ {n} ^ {\prime} = \rho (\mathrm{C} ^ {\prime})\right)
\]

of smallest possible length such that \(F \subset C_{0}^{\prime}\) (consider an apartment \(A^{\prime}\) containing \(C^{\prime}\) and F and apply b) and the fact that the restriction of \(\rho\) to \(A^{\prime}\) is an isomorphism).

\begin{enumerate}
\def\labelenumi{\alph{enumi})}
\setcounter{enumi}{3}
\tightlist
\item
  Let A be an apartment, \(C_{1}\) and \(C_{2}\) two distinct neighbouring chambers contained in A, \(C'\) a chamber containing \(C_{1} \cap C_{2}\) and distinct from \(C_{1}\) and \(C_{2}\) , and \(A_{i}'\) an apartment containing \(C_{i}\) and \(C'\) (for i = 1, 2). Let \(\varphi_{i}\) (resp. \(\psi_{i}\) ) be the retraction of I onto A (resp. \(A_{i}'\) ) with centre \(C_{i}\) (resp. \(C'\) ) and let \(\rho_{i}\) be the restriction of \(\varphi_{i} \circ \psi_{i}\) to A. Let C be a chamber of A; show that if
\end{enumerate}

\[
d (\mathrm{C}, \mathrm{C} _ {1}) \leqslant d (\mathrm{C}, \mathrm{C} _ {2}),
\]

then \(\rho_{1}(a) = a\) for all \(a \in \mathbb{C}\) and \(d(\mathbb{C}, \mathbb{C}_{1}) < d(\mathbb{C}, \mathbb{C}_{2})\) (consider a minimal gallery \(\Gamma\) with extremities \(\mathbb{C}\) and \(\mathbb{C}_{1}\) and apply \(c\) ) to show that \(\rho_{1}(\Gamma)\) is minimal; then use Exerc. 15 b)). Show that, if \(d(\mathbb{C}, \mathbb{C}_{2}) < d(\mathbb{C}, \mathbb{C}_{1})\) , then \(\rho_{1}(\mathbb{C}) \neq \mathbb{C}\) and \(\rho_{1}^{2}(\mathbb{C}) = \rho_{1}(\mathbb{C})\) (take a minimal gallery \(\Gamma\) with extremities \(\mathbb{C}\) and \(\mathbb{C}_{2}\) and use \(c\) ) to show that \(\rho_{1}(\Gamma)\) is a gallery of smallest possible length having one extremity equal to \(\rho_{1}(\mathbb{C})\) and the other containing \(\mathbb{C}_{1} \cap \mathbb{C}_{2}\) ; deduce that \(d(\mathbb{C}_{1}, \rho_{1}(\mathbb{C}_{1})) \leqslant d(\mathbb{C}_{2}, \rho_{1}(\mathbb{C}_{1}))\) .

Show that \(\rho_{i}\) is a folding (Exerc. 19) of A (show that \(A = \rho_{1}(A) \cup \rho_{2}(A)\) and define an involutive automorphism \(\sigma\) of A by putting \(\sigma(a) = \rho_{2}(a)\) if \(a \in \rho_{1}(A)\) and \(\sigma(a) = \rho_{1}(a)\) if \(a \in \rho_{2}(A)\) ; show that if C is a chamber contained in \(\rho_{1}(A)\) , then \(\rho_{1}^{-1}(C) = \{C, \rho_{2}(C)\}\) ).

\begin{enumerate}
\def\labelenumi{\alph{enumi})}
\setcounter{enumi}{4}
\tightlist
\item
  Let I be a spacious structured building, that is to say that every panel is contained in at least three chambers. Show that there exists a Coxeter system (W,S), unique up to isomorphism, such that the apartments of I are isomorphic to the apartment \(A_{0}\) associated to (W,S) (use d) and Exerc. 19). Let A be an apartment of I and \(\varphi\) an isomorphism from \(A_{0}\) to A. Show that there exists a unique numbering of I, with values in S, such that the types of a and \(\varphi(a)\) are equal for all \(a \in A_{0}\) (choose a chamber C of A and show, by using b), that, if \(A'\) and \(A''\) are two apartments of I containing C, the numberings of \(A'\) and \(A''\) extending that determined on C coincide on \(A' \cap A''\) ).
\end{enumerate}

We shall say that the Coxeter system (W,S) and the numbering thus obtained are adapted to I.

Show that the retractions introduced in a) are allowed endomorphisms. Show that a subset of I contained in an apartment A of I is convex in I if and only if it is convex in A (Exerc. 20).

\(f\) ) We retain the notation of \(e\) ) and let \(\mathfrak{A}\) be the set of allowed isomorphisms from \(\mathbf{A}_0\) to the various apartments of I. Show that, if \(\varphi, \psi \in \mathfrak{A}\) and if C is a chamber and F a facet of I contained in \(\varphi(\mathbf{A}_0) \cap \psi(\mathbf{A}_0)\) , there exists an element \(w \in W\) such that \(\psi^{-1}(C) = w\varphi^{-1}(C)\) and \(\psi^{-1}(F) = w\varphi^{-1}(F)\) (consider an isomorphism \(\lambda\) from \(\varphi(\mathbf{A}_0)\) to \(\psi(\mathbf{A}_0)\) leaving fixed the points of \(\varphi(\mathbf{A}_0) \cap \psi(\mathbf{A}_0)\) and apply Exerc. 21 \(b\) ) to the automorphism \(\psi^{-1} \circ \lambda \circ \varphi^{-1}\) of \(\mathbf{A}_0\) ).

\begin{enumerate}
\def\labelenumi{\arabic{enumi})}
\setcounter{enumi}{24}
\tightlist
\item
  Let I be a set and let \(\mathfrak{F}\) be the set of finite subsets of I. For any \(A \in \mathfrak{F}\) , put \(\varepsilon(A) = (-1)^{\operatorname{Card}(A)}\) . Let \(G\) be an abelian group written additively, and let \(\varphi\) and \(\psi\) be two maps from \(\mathfrak{F}\) to \(G\) . Show that the following properties are equivalent:
\end{enumerate}

\[
\varphi (\mathrm{A}) = \sum_ {\mathrm{B} \subset \mathrm{A}} \psi (\mathrm{B}) \quad \text {   for   all   } \mathrm{A} \in \mathfrak {F};\tag{i}
\]

\[
\psi (A) = \sum_ {B \subset A} \varepsilon (A - B) \varphi (B) \quad \text {   for   all   } A \in \mathfrak {F}.\tag{ii}
\]

\begin{enumerate}
\def\labelenumi{\arabic{enumi})}
\setcounter{enumi}{25}
\tightlist
\item
  Let \((W, S)\) be a Coxeter system, with S finite. For any subset H of W, denote by \(\mathrm{H}(t)\) the formal power series with integer coefficients defined by
\end{enumerate}

\[
\mathrm{H} (t) = \sum_ {w \in \mathrm{H}} t ^ {l (w)}.
\]

\begin{enumerate}
\def\labelenumi{\alph{enumi})}
\tightlist
\item
  Assume that \(\operatorname{Card}(S) = 2\) . Show that
\end{enumerate}

\[
\mathrm{W} (t) = \left(1 + t - t ^ {m} - t ^ {m + 1}\right) / (1 - t) \quad \text { if   } \mathrm{W} \text {   is   of   finite   order   } 2 m,
\]

\[
\mathrm{W} (t) = (1 + t) / (1 - t) \quad \text { if   } \mathrm{W} \text {   is   infinite. }
\]

\begin{enumerate}
\def\labelenumi{\alph{enumi})}
\setcounter{enumi}{1}
\tightlist
\item
  Assume that W is finite. Let \(w_0\) be the element of W of greatest length (Exerc. 22) and let \(m = l(w_0)\) . Show that
\end{enumerate}

\[
\mathrm{W} (t) = t ^ {m} \mathrm{W} (t ^ {- 1})
\]

(use Exerc. 22).

\begin{enumerate}
\def\labelenumi{\alph{enumi})}
\setcounter{enumi}{2}
\tightlist
\item
  Let X be a subset of S. Denote by \(A_{X}\) the set of \((X,\varnothing)\) -reduced elements of W (Exerc. 3) and \(W_{X}\) the subgroup of W generated by X. We know (Exerc. 3) that an element \(w \in W\) belongs to \(A_{X}\) if and only if \(l(xw) = l(w) + 1\) for all \(x \in X\) , that every element \(w \in W\) can be written uniquely in the form w = uv with \(u \in W_{X}\) and \(v \in A_{X}\) , and that in that case \(l(w) = l(u) + l(v)\) . Deduce the formula
\end{enumerate}

\[
\mathrm{W} (t) = \mathrm{W} _ {\mathrm{X}} (t). \mathrm{A} _ {\mathrm{X}} (t).
\]

\(d\) ) Retain the above notation and denote by \(B_{X}\) the set of \(w\in A_X\) such that \(l(sw) = l(w) - 1\) for all \(s\in S - X\) . Show that

\[
A _ {X} (t) = \sum_ {X \subset Y \subset S} B _ {Y} (t).
\]

Deduce that

\[
\mathrm{B} _ {\mathrm{X}} (t) = \sum_ {\mathrm{X} \subset \mathrm{Y} \subset \mathrm{S}} \varepsilon (\mathrm{Y} - \mathrm{X}) \mathrm{A} _ {\mathrm{Y}} (t) \quad \text { with } \quad \varepsilon (\mathrm{Z}) = (- 1) ^ {\operatorname{Card} (\mathrm{Z})}
\]

(use Exerc. 25).

\begin{enumerate}
\def\labelenumi{\alph{enumi})}
\setcounter{enumi}{4}
\tightlist
\item
  Assume that W is finite and define m and \(w_{0}\) as in b). Show that \(B_{\varnothing} = \{w_{0}\}\) and that
\end{enumerate}

\[
t ^ {m} = \sum_ {\mathrm{Y} \subset \mathrm{S}} \varepsilon (\mathrm{Y}) \frac {\mathrm{W} (t)}{\mathrm{W} _ {\mathrm{Y}} (t)}
\]

(use c) and d)).

\begin{enumerate}
\def\labelenumi{\alph{enumi})}
\setcounter{enumi}{5}
\tightlist
\item
  Assume that W is infinite. Show that \(B_{\varnothing} = \varnothing\) and that
\end{enumerate}

\[
0 = \sum_ {Y \subset S} \frac {\varepsilon (Y)}{W _ {Y} (t)}.
\]

\(g\) ) Show that the formal power series \(W(t)\) is a rational function of \(t\) (use \(f\) ) and argue by induction on \(\operatorname{Card}(\mathbf{S})\) ). Show that this rational function vanishes whenever \(t\) is a root of unity, and that \(1 / W(\infty)\) is an integer. Show that \(\frac{1}{W(t^{-1})} \in \mathbf{Z}[[t]]\) .

